% TOP_PROBLEM: {"id": 10000046, "title": "Nonintersecting couplings of random walks in dimensions three and four", "category_id": 19, "category": {"id": 19, "name": "probability", "display_name": "Probability", "description": "Problems involving probability theory, stochastic processes, and random structures.", "slug": "probability", "order_index": 19, "created_at": "2026-07-31T15:26:25.671Z"}, "status": "open", "problem_number": "AMR-099-0046", "set_id": 13}
% FIRST_POSTED: 2026-10-01
% ENTRY_KIND: statement_only
% UnsolvedMath ID 10000046; source catalogue code AMR-099-0046
% !TeX program = lualatex
% Definition and sources only; this file is not an LLM solution attempt.
\documentclass[11pt,a4paper]{article}
\usepackage[margin=25mm]{geometry}
\usepackage{fontspec}
\setmainfont{lmroman10-regular.otf}[BoldFont=lmroman10-bold.otf,
 ItalicFont=lmroman10-italic.otf,BoldItalicFont=lmroman10-bolditalic.otf]
\usepackage{amsmath,amssymb,mathtools}
\usepackage{unicode-math}
\setmathfont{latinmodern-math.otf}
\AtBeginDocument{\let\setminus\smallsetminus}
\usepackage{xurl,hyperref}
\hypersetup{colorlinks=true,urlcolor=blue}
\setcounter{secnumdepth}{0}
\setlength{\parindent}{0pt}
\setlength{\parskip}{5pt}
\setlength{\emergencystretch}{3em}
\begin{document}
\section*{Nonintersecting couplings of random walks in dimensions three and four}\label{10000046}
{\small UnsolvedMath ID 10000046 · AMR-099-0046 · Probability · AMR Open Problem Lists}
\subsection{Definitions and mathematical statement}
For $d\in\{3,4\}$, give $\mathbb Z^d$ its nearest-neighbor graph structure. Fix vertices $x,y$ at graph distance $\|x-y\|_1=10$. Let $P_x$ and $P_y$ denote the probability laws on $(\mathbb Z^d)^{\mathbb N_0}$ of simple random walks started at these vertices: at each step, independently of the preceding increments, the walk chooses one of the $2d$ coordinate neighbors with equal probability.

A coupling is any probability law $Q$ on pairs of paths $(X,Y)$ whose first marginal is $P_x$ and whose second marginal is $P_y$. Does there exist such a coupling for which
\[
Q\bigl(\{X_n:n\ge0\}\cap\{Y_m:m\ge0\}=\varnothing\bigr)>0?
\]
Ask this for each of the two dimensions and every pair of starting points at the specified distance.

The disjointness condition forbids intersections at different times as well as simultaneous collisions. Each entire marginal trajectory must still have the ordinary simple-random-walk law, but the two trajectories may depend on one another arbitrarily. In particular, the problem does not impose a Markovian or co-adapted coupling, under which both walks would need to keep their transition rules relative to the joint past. The successful disjointness event need only have positive probability; no uniform lower bound over starting pairs is requested. Translating both starts together does not change the question, while graph distance ten permits several relative directions that are not related by lattice symmetries.
\subsection{Short English statement}
Can two ordinary random walks in three or four dimensions be coupled so their entire ranges are disjoint with positive probability?
\subsection{Sources}
\begin{itemize}
\item[S1] ULAM.AI and the UnsolvedMath Contributors, supplied problems.json, record 10000046 (AMR-099-0046), AMR Open Problem Lists. Catalogue identity and source transcription; CC BY 4.0.
\url{https://huggingface.co/datasets/ulamai/UnsolvedMath}.
\item[S2] Benjamini - Coarse Geometry and Randomness (Saint-Flour notes), Open problem 12.33, PDF page 104. Original source indexed by AMR-099-0046.
\url{https://www.wisdom.weizmann.ac.il/~itai/stflouraug24.pdf#page=104}.
\item[S3] I. Benjamini, Coarse Geometry and Randomness, primary Saint-Flour notes; the numbered question and its surrounding definitions.
\url{https://arquivo.pt/noFrame/replay/20201231041548id_/http://www.wisdom.weizmann.ac.il/~itai/stflouraug24.pdf}.
\end{itemize}
\end{document}
